% GENERATED FILE. Edit canonical lessons, then run npm run generate:books.
% canonical-source-hash: fnv1a64:7dd596b6914ed714
% canonical-lessons: TE-C431-pramanam, TE-C431-sthayi, TE-C431-dhorani, TE-C431-samarthavantamaina, TE-C431-pariksincu

\chapter{Standards and Levels}
\label{ch:te-standards-and-levels}
\hlchaptermodality{431}

\begin{chapteropening}
\textbf{By the end of this chapter:} \emph{I can talk about standards, levels, trends and testing.}

Complete the last lesson of chapter 431: I can talk about standards, levels, trends and testing.
\end{chapteropening}

\section[\texorpdfstring{\te{ప్రమాణం} (pramāṇaṁ)}{ప్రమాణం (pramāṇaṁ)}]{\te{ప్రమాణం} (pramāṇaṁ) — a standard; a norm}

\label{lesson:TE-C431-pramanam}



\begin{quote}
Before the new one: say the Telugu for to control, then the Telugu for a limit.
\end{quote}

\subsection*{You'll want to know: \te{ప్రమాణం}}
\textbf{\te{ప్రమాణం}} — \emph{pramāṇaṁ} — \textquotedblleft{}a standard; a norm\textquotedblright{}.

\begin{quote}
\te{ఆ} \te{కళాశాలలో} \te{విద్యా} \te{ప్రమాణాలు} \te{ఎక్కువ}.
\end{quote}

\begin{quote}
\emph{ā kaḷāśālalō vidyā pramāṇālu ekkuva.}
\end{quote}

\begin{quote}
Academic standards are high at that college.
\end{quote}

Plural: \te{ప్రమాణాలు}, \emph{pramāṇālu}.

Also an oath: \te{ప్రమాణం} \te{చేయు}, to swear.

\begin{cousinweb}
\textbf{\te{ప్రమాణం}} (\emph{pramāṇaṁ}) is borrowed from Sanskrit \emph{pramāṇa}, \textquotedblleft{}measure; a standard, an authority; proof\textquotedblright{} (CDIAL 8723).
\end{cousinweb}

\begin{grammarlens}[title={What you've built}]
One more word for saying what happens if.
\end{grammarlens}

\subsection*{Your turn}
\begin{itemize}
\raggedright
  \item \emph{Say it:} \emph{pramāṇaṁ}
  \item \emph{Say it:} \emph{pramāṇaṁ}, once more
  \item \emph{Say it:} \emph{ā kaḷāśālalō vidyā pramāṇālu ekkuva.} — Academic standards are high at that college.
  \item \emph{Recall:} say the Telugu for to control, then the Telugu for a limit, then say \emph{pramāṇaṁ} again
\end{itemize}

\begin{tcolorbox}[breakable,colback=teal!4,colframe=teal!35!black,title={Before you move on}]
What does \te{ప్రమాణం} mean? (\textquotedblleft{}A standard; a norm\textquotedblright{}.) Say it once more.
\end{tcolorbox}
\section[\texorpdfstring{\te{స్థాయి} (sthāyi)}{స్థాయి (sthāyi)}]{\te{స్థాయి} (sthāyi) — a level; a standing}

\label{lesson:TE-C431-sthayi}



\begin{quote}
Before the new one: say the Telugu for a limit, then the Telugu for a standard.
\end{quote}

\subsection*{You'll want to know: \te{స్థాయి}}
\textbf{\te{స్థాయి}} — \emph{sthāyi} — \textquotedblleft{}a level; a standing\textquotedblright{}.

\begin{quote}
\te{ఆమె} \te{తెలుగు} \te{ఇప్పుడు} \te{మంచి} \te{స్థాయిలో} \te{ఉంది}.
\end{quote}

\begin{quote}
\emph{āme telugu ippuḍu mañci sthāyilō undi.}
\end{quote}

\begin{quote}
Her Telugu is at a good level now.
\end{quote}

\te{జాతీయ} \te{స్థాయిలో}, at the national level.

\begin{cousinweb}
\textbf{\te{స్థాయి}} (\emph{sthāyi}) is borrowed from Sanskrit \emph{sthāyin}, \textquotedblleft{}standing\textquotedblright{}, and \textquotedblleft{}lasting\textquotedblright{} (CDIAL 13764), from \emph{sthā}, \textquotedblleft{}to stand\textquotedblright{}.
\end{cousinweb}

\begin{grammarlens}[title={What you've built}]
One more word for saying what happens if.
\end{grammarlens}

\subsection*{Your turn}
\begin{itemize}
\raggedright
  \item \emph{Say it:} \emph{sthāyi}
  \item \emph{Say it:} \emph{sthāyi}, once more
  \item \emph{Say it:} \emph{āme telugu ippuḍu mañci sthāyilō undi.} — Her Telugu is at a good level now.
  \item \emph{Recall:} say the Telugu for a limit, then the Telugu for a standard, then say \emph{sthāyi} again
\end{itemize}

\begin{tcolorbox}[breakable,colback=teal!4,colframe=teal!35!black,title={Before you move on}]
What does \te{స్థాయి} mean? (\textquotedblleft{}A level; a standing\textquotedblright{}.) Say it once more.
\end{tcolorbox}
\section[\texorpdfstring{\te{ధోరణి} (dhōraṇi)}{ధోరణి (dhōraṇi)}]{\te{ధోరణి} (dhōraṇi) — a trend; a tendency}

\label{lesson:TE-C431-dhorani}



\begin{quote}
Before the new one: say the Telugu for a standard, then the Telugu for a level.
\end{quote}

\subsection*{You'll want to know: \te{ధోరణి}}
\textbf{\te{ధోరణి}} — \emph{dhōraṇi} — \textquotedblleft{}a trend; a tendency\textquotedblright{}.

\begin{quote}
\te{యువతలో} \te{ఒక} \te{కొత్త} \te{ధోరణి} \te{కనిపిస్తోంది}.
\end{quote}

\begin{quote}
\emph{yuvatalō oka kotta dhōraṇi kanipistōndi.}
\end{quote}

\begin{quote}
A new trend can be seen among young people.
\end{quote}

Also a person’s manner: \te{అతని} \te{ధోరణి} \te{నాకు} \te{నచ్చలేదు}, I did not like his attitude.

\begin{cousinweb}
\textbf{\te{ధోరణి}} (\emph{dhōraṇi}) may go back to Sanskrit \emph{dhoraṇī}, \textquotedblleft{}an uninterrupted series\textquotedblright{} (Monier-Williams), but the link is not settled.
\end{cousinweb}

\begin{grammarlens}[title={What you've built}]
One more word for saying what happens if.
\end{grammarlens}

\subsection*{Your turn}
\begin{itemize}
\raggedright
  \item \emph{Say it:} \emph{dhōraṇi}
  \item \emph{Say it:} \emph{dhōraṇi}, once more
  \item \emph{Say it:} \emph{yuvatalō oka kotta dhōraṇi kanipistōndi.} — A new trend can be seen among young people.
  \item \emph{Recall:} say the Telugu for a standard, then the Telugu for a level, then say \emph{dhōraṇi} again
\end{itemize}

\begin{tcolorbox}[breakable,colback=teal!4,colframe=teal!35!black,title={Before you move on}]
What does \te{ధోరణి} mean? (\textquotedblleft{}A trend; a tendency\textquotedblright{}.) Say it once more.
\end{tcolorbox}
\section[\texorpdfstring{\te{సమర్థవంతమైన} (samarthavantamaina)}{సమర్థవంతమైన (samarthavantamaina)}]{\te{సమర్థవంతమైన} (samarthavantamaina) — effective; efficient}

\label{lesson:TE-C431-samarthavantamaina}



\begin{quote}
Before the new one: say the Telugu for a level, then the Telugu for a trend.
\end{quote}

\subsection*{You'll want to know: \te{సమర్థవంతమైన}}
\textbf{\te{సమర్థవంతమైన}} — \emph{samarthavantamaina} — \textquotedblleft{}effective; efficient\textquotedblright{}.

\begin{quote}
\te{ఆమె} \te{సమర్థవంతమైన} \te{అధికారి}.
\end{quote}

\begin{quote}
\emph{āme samarthavantamaina adhikāri.}
\end{quote}

\begin{quote}
She is an efficient officer.
\end{quote}

\te{సమర్థ}, capable, + -\te{వంత}, having, + -\te{మైన}. \te{సమర్థవంతంగా}, effectively.

\begin{cousinweb}
\textbf{\te{సమర్థవంతమైన}} (\emph{samarthavantamaina}) is built on Sanskrit \emph{samartha}, \textquotedblleft{}appropriate; capable\textquotedblright{} (CDIAL 13189), with \emph{-vant}, \textquotedblleft{}having\textquotedblright{}, and Telugu -\te{మైన} for an adjective.
\end{cousinweb}

\begin{grammarlens}[title={What you've built}]
One more word for saying what happens if.
\end{grammarlens}

\subsection*{Your turn}
\begin{itemize}
\raggedright
  \item \emph{Say it:} \emph{samarthavantamaina}
  \item \emph{Say it:} \emph{samarthavantamaina}, once more
  \item \emph{Say it:} \emph{āme samarthavantamaina adhikāri.} — She is an efficient officer.
  \item \emph{Recall:} say the Telugu for a level, then the Telugu for a trend, then say \emph{samarthavantamaina} again
\end{itemize}

\begin{tcolorbox}[breakable,colback=teal!4,colframe=teal!35!black,title={Before you move on}]
What does \te{సమర్థవంతమైన} mean? (\textquotedblleft{}Effective; efficient\textquotedblright{}.) Say it once more.
\end{tcolorbox}
\section[\texorpdfstring{\te{పరీక్షించు} (parīkṣiñcu)}{పరీక్షించు (parīkṣiñcu)}]{\te{పరీక్షించు} (parīkṣiñcu) — to test out; to examine}

\label{lesson:TE-C431-pariksincu}



\begin{quote}
Before the new one: say the Telugu for a trend, then the Telugu for effective.
\end{quote}

\subsection*{You'll want to know: \te{పరీక్షించు}}
\textbf{\te{పరీక్షించు}} — \emph{parīkṣiñcu} — \textquotedblleft{}to test out; to examine\textquotedblright{}.

\begin{quote}
\te{వైద్యుడు} \te{నా} \te{కళ్ళను} \te{పరీక్షించాడు}.
\end{quote}

\begin{quote}
\emph{vaidyuḍu nā kaḷḷanu parīkṣiñcāḍu.}
\end{quote}

\begin{quote}
The doctor examined my eyes.
\end{quote}

Past (I …): \te{పరీక్షించాను}, \emph{parīkṣiñcānu}.

The verb of \te{పరీక్ష}, an exam, a test, which you know.

\begin{cousinweb}
\textbf{\te{పరీక్షించు}} (\emph{parīkṣiñcu}) is built on Sanskrit \emph{parīkṣā}, \textquotedblleft{}inspection, investigation, examination, test\textquotedblright{} (Monier-Williams), with the Telugu verb ending -\te{ఇంచు}; \te{పరీక్ష} is the same word.
\end{cousinweb}

\begin{grammarlens}[title={What you've built}]
One more word for saying what happens if.
\end{grammarlens}

\subsection*{Your turn}
\begin{itemize}
\raggedright
  \item \emph{Say it:} \emph{parīkṣiñcu}
  \item \emph{Say it:} \emph{parīkṣiñcu}, once more
  \item \emph{Say it:} \emph{vaidyuḍu nā kaḷḷanu parīkṣiñcāḍu.} — The doctor examined my eyes.
  \item \emph{Recall:} say the Telugu for a trend, then the Telugu for effective, then say \emph{parīkṣiñcu} again
\end{itemize}

\begin{tcolorbox}[breakable,colback=teal!4,colframe=teal!35!black,title={Before you move on}]
What does \te{పరీక్షించు} mean? (\textquotedblleft{}To test out; to examine\textquotedblright{}.) Say it once more.
\end{tcolorbox}
